\documentclass[11pt]{article}
\usepackage{fontspec}
\setmainfont{Times New Roman}
\setsansfont{Arial}
\setmonofont{Consolas}
\usepackage{amsmath,amssymb}
\usepackage{geometry}
\usepackage{microtype}
\geometry{a4paper,margin=3cm}
\setlength{\parindent}{1.25em}
\setlength{\parskip}{0pt}
\emergencystretch=6em
\allowdisplaybreaks

\begin{document}

\section*{\S\ 1. Означеньа и определьеньа. Свод познаных резултатов.}

Како обычно, означајемо ряд пременных $x$ како ряд елементов $x_1,x_2,\ldots,x_n$, и аналогично ряд пременных $p_\rho$ ($\rho=1,2,\ldots,n-1$) како ряд из $\binom n\rho$ елементов $p_{i_1i_2\ldots i_\rho}$, кде $p_{i_1i_2\ldots i_\rho}$ означује $\binom n\rho$ детерминантов, утворених из матрице $\rho$ рядов $x^{(1)},x^{(2)},\ldots,x^{(\rho)}$:
\[
\left|\begin{array}{cccc}
 x_{i_1}^{(1)}&x_{i_2}^{(1)}&\cdots&x_{i_\rho}^{(1)}\\
 x_{i_1}^{(2)}&x_{i_2}^{(2)}&\cdots&x_{i_\rho}^{(2)}\\
 \vdots&\vdots&&\vdots\\
 x_{i_1}^{(\rho)}&x_{i_2}^{(\rho)}&\cdots&x_{i_\rho}^{(\rho)}
\end{array}\right|,
\qquad (i_1<i_2<\cdots<i_\rho=1,2,\ldots,n).
\]
По теореме о одговарјајучих матрицах\footnote{Срав., на примјер: Gordan--Kerschensteiner, \emph{Vorlesungen über Invariantentheorie}, Бд. I, С. 99.} всакому реду $p_\rho$ взаимно једнозначно приряджен јест други ряд $q_{n-\rho}$ чрез релацију, која плати за всак елемент реда:
\begin{equation}
 p_{i_1i_2\ldots i_\rho}=(-1)^{\frac{\rho(\rho+1)}{2}+i_1+i_2+\cdots+i_\rho}q_{j_1j_2\ldots j_{n-\rho}},
\tag{1}
\end{equation}
кде $i$ и $j$ сут одделно упоряджене и вместе творјат полну пермутацију чисел од 1 до $n$. Ряд $q_{n-\rho}$ состаји из $\binom n\rho$ детерминантов ступньа $n-\rho$, $q_{j_1j_2\ldots j_{n-\rho}}$, утворених из матрице $n-\rho$ рядов $u^{(1)},u^{(2)},\ldots,u^{(n-\rho)}$, контрагредиентных к $x$. Те ряды задовольајут $\rho(n-\rho)$ условным уравненьам:
\begin{equation}
 (u^{(k)}\mid x^{(l)})=\sum_{\alpha=1}^{n}u_\alpha^{(k)}x_\alpha^{(l)}=0;
 \qquad (k=1,2,\ldots,n-\rho;\ l=1,2,\ldots,\rho).
\tag{2}
\end{equation}
С Цлебсцхом\footnote{а. а. О. \S\ 5.} будемо јих сматрати тако нормованими, же фактор пропорционалности, кторы бы иначе входил в (1.), поставльен јест равным јединици.

По Цлебсцху\footnote{а. а. О. \S\ 12.} (и такоже по позднејших изследованьах Ф. Деруытса и А. Цапеллија, срав. Ввод), најобче формы можно свести к таким, кторе обсахујут највише једен из всакого од $n-1$ рядов $p_\rho$, то јест симболично се представјајут тако:
\begin{equation}
 (S^1\mid p_1)^{m_1}(S^2\mid p_2)^{m_2}\cdots(S^{n-1}\mid p_{n-1})^{m_{n-1}},
\tag{3}
\end{equation}
кде, по аналогији с (2.), поставльено јест
\[
 (S^\rho\mid p_\rho)=\sum S^{i_1i_2\ldots i_\rho}p_{i_1i_2\ldots i_\rho}.
\]
Из-за нелинеарных идентичных релациј меджу елементами реда $p_\rho$ ряды симболов $S^\rho$ можно нормовати тако, же оны задовольајут тым самым идентичностјам, то јест ведут се како детерминанты $\rho$-рядној матрице, ктореј ряды $s^{(1)},s^{(2)},\ldots,s^{(\rho)}$ сут контрагредиентне к $x$.\footnote{Clebsch, а. а. О. \S\ 4.} Зато всакому реду $S^\rho$ можно взаимно једнозначно прирядити други ряд $S_{n-\rho}$ чрез релацију:
\begin{equation}
 S_{i_1i_2\ldots i_{n-\rho}}=(-1)^{\frac{(n-\rho)(n-\rho+1)}2+i_1+i_2+\cdots+i_{n-\rho}}S^{j_1j_2\ldots j_\rho},
\tag{4}
\end{equation}
кде знову $i$ и $j$ сут одделно упоряджене и вместе точно изчерпавајут числа од 1 до $n$. Елементы $S_{i_1i_2\ldots i_{n-\rho}}$ реда $S_{n-\rho}$ ведут се како детерминанты, утворене из $(n-\rho)$ рядов $\sigma^{(1)},\sigma^{(2)},\ldots,\sigma^{(n-\rho)}$, когредиентных к $x$, а ряды $s$ и $\sigma$ сут связане чрез $\rho(n-\rho)$ условиј:
\begin{equation}
 (s^{(k)}\mid \sigma^{(l)})=0,\qquad(k=1,2,\ldots,\rho;\ l=1,2,\ldots,n-\rho).
\tag{5}
\end{equation}
Из релациј (1.) и (4.) следује, же всакы фактор в (3.) допушча еквивалентно четверократно симболично представјенье:
\begin{equation}
\begin{aligned}
 (S^\rho\mid p_\rho)&=(S^\rho q_{n-\rho})=(S_{n-\rho}p_\rho)\\
 &=(-1)^{\rho(n-\rho)}(q_{n-\rho}\mid S_{n-\rho}),
\end{aligned}
\tag{6}
\end{equation}
кде, како всегда далье, $(S^\rho q_{n-\rho})$ и $(S_{n-\rho}p_\rho)$ сут скраченьа за детерминанты $(s^{(1)}\cdots s^{(\rho)}u^{(1)}\cdots u^{(n-\rho)})$ и $(\sigma^{(1)}\cdots\sigma^{(n-\rho)}x^{(1)}\cdots x^{(\rho)})$; чрез Лаплацеово разложенье показује се, же оны зависят само од рядов $S^\rho,q_{n-\rho}$ односно $S_{n-\rho},p_\rho$, што имаје означити выше дано писанье. Изразы $(S^\rho\mid p_\rho)$ и $(q_{n-\rho}\mid S_{n-\rho})$ означајут, по выше даном објашньеньу, продукты матриц рядов $s$ и $x$, односно рядов $u$ и $\sigma$, при чем под матричным продуктом разумеје се сума продуктов одговарјајучих детерминантов, то јест вредност детерминанта продуктној матрице.

Карактеристично число $\rho(n-\rho)$, кторе принадлежи реду симболов (односно реду пременных), называмо по Цлебсцху вагоју реда;\footnote{а. а. О. \S\ 11.} число $\rho$, кторе дава число рядов $s$, горным индексом вагы, а число $n-\rho$, кторе дава число рядов $\sigma$, долным индексом вагы. Из релациј (4.) следује, же ряд симболов може наступити в једнеј и тој самој релацији једен раз с горным, а други раз с долным индексом вагы; тоже плати за наступанье одповедајучих рядов $p_\rho$ и $q_{n-\rho}$. Јешче замечајемо, же обчно означајемо: ряды пременных вышего нивоја, кторых одговарјајуча матрица состаји из рядов $x$, чрез $p$; те, кторых матрица состаји из рядов $u$, чрез $q$; ряды симболов, когредиентне к $x$, малыми гречскыми буквами $\sigma,\tau,\alpha,\beta$; ряды симболов, когредиентне к $u$, малыми латинскыми буквами $s,t,a,b$; все друге ряды симболов великыми латинскыми буквами с индексом вагы: $S^\rho$, $T^\tau$, $A^\rho$ односно $S_{n-\sigma}$, $T_{n-\tau}$, $A_{n-\rho}$. Согласно горному або долному индексу вагы пишемо такоже поједине елементы реда с горными або долными индексами, како на примјер в (4.).

\end{document}

